% Source-derived chapter 01: Introduction
% Original source: tates-thesis-de-rham-setting
\chapter{Introduction}
\label{tates-thesis-de-rham-setting-chapter-01}
\source{tates-thesis-de-rham-setting}

\begin{remark}[Source section]
\label{tates-thesis-de-rham-setting-chapter-01-source}
\source{tates-thesis-de-rham-setting}
\source{tates-thesis-de-rham-setting}
This chapter reproduces the corresponding section of the retrieved
LaTeX source, including its definitions, statements, examples, and
proofs. It has not yet been translated into Lean.
\notready
\end{remark}

% The source section title was: Introduction
\label{s:intro}
\source{tates-thesis-de-rham-setting}

\subsection{Statement of the main results}

We work over a field $k$ of characteristic zero. 

Let $\sY$ be the moduli space of rank $1$ de Rham local systems on 
the punctured disc equipped with a flat section, and 
let $\fL \bA^1$ denote the algebraic loop space of $\bA^1$.

Our main result asserts:

\begin{thm}[Thm. \ref{t:main-precise}]
\source{tates-thesis-de-rham-setting}
\notready\label{t:main}
\source{tates-thesis-de-rham-setting}

There is a canonical equivalence of DG categories:
%
\[
\Delta:D^!(\fL \bA^1) \simeq \IndCoh^*(\sY).
\]

Moreover, this equivalence is compatible with the local geometric 
class field theory of Beilinson-Drinfeld. 

\end{thm}

Here the left hand side is the category of $D$-modules on $\fL \bA^1$,
as defined in \cite{dario-*/!}, \cite{dmod}. The right hand side
is the category of ind-coherent sheaves on $\sY$, which we construct
in \S \ref{s:indcoh} (following \cite{methods}). In both cases,
there are infinite-dimensional subtleties in the definitions; 
these are far more severe on the right hand side. 

\begin{rem}
\source{tates-thesis-de-rham-setting}
\notready

The assertions of the theorem are made more precise in the body of the paper. 
For now, we content ourselves with the following remark on local class field theory.

Beilinson-Drinfeld construct an equivalence:
%
\begin{equation}\label{eq:intro-lcft}
\source{tates-thesis-de-rham-setting}
D^*(\fL\bG_m) \simeq \QCoh(\LocSys_{\bG_m})
\end{equation}

\noindent for $\LocSys_{\bG_m}$ the moduli space of rank $1$ de Rham local systems on the
punctured disc. The left hand side here naturally acts on 
the left hand side of Theorem \ref{t:main}, while the right hand side
naturally acts on the right hand side of Theorem \ref{t:main}. The compatibility
asserts that the equivalence of Theorem \ref{t:main} is compatible with 
these actions.

\end{rem}

\begin{rem}
\source{tates-thesis-de-rham-setting}
\notready

Both sides of the above equivalence have natural $t$-structures.
However, this equivalence is \emph{not} $t$-exact; it is necessarily
an equivalence of derived categories. This is in contrast to \eqref{eq:intro-lcft},
which largely amounts to an equivalence of abelian categories. 

\end{rem}

\begin{rem}
\source{tates-thesis-de-rham-setting}
\notready

For $K$ a local field, Tate's thesis \cite{tate-thesis} 
(see also \cite{tate-thesis-weil}) 
considers the decomposition of the space
$\sD(K)$ of tempered distributions on $K$ 
as a representation for $K^{\times}$. We consider Theorem \ref{t:main} 
as an analogue of the arithmetic situation. We observe that
our result applies not only at the level of eigenspaces, 
but describes a categorical analogue of the whole space $\sD(K)$.

We plan to return to global aspects of the subject in future work. 

%% reference?

\end{rem}

\begin{rem}
\source{tates-thesis-de-rham-setting}
\notready

As far as we are aware, our work is the first one in geometric representation
theory to prove a theorem about coherent sheaves on a 
space of the form $\Maps(\o{\sD}_{dR},Y)$ for $Y$ an Artin stack
that is not a scheme and is not a classifying stack (for us: $Y = \bA^1/\bG_m$). 
As we will see in \S \ref{s:spectral} and \ref{s:indcoh}, there are substantial
technical difficulties that arise in this setting.
Roughly speaking, the singularities of the space $\sY$ are genuinely of infinite type. 
We can only overcome these difficulties in our specific (abelian) setting.

\end{rem}

\begin{rem}
\source{tates-thesis-de-rham-setting}
\notready

As discussed later in \S \ref{ss:intro-strategy}, the main piece of our
construction is a realization of $D$-modules on $\bA^n$ in spectral terms;
this is the content of \S \ref{s:weyl}.
Already in the $n = 1$ case, our realization of $D$-modules on 
$\bA^1$ as coherent sheaves on some space is novel; in this case, 
we highlight that the expression
is as coherent sheaves on the subspace $\sY^{\leq 1} \subset \sY$ 
discussed in detail in  \S \ref{ss:rs-sketch}.

\end{rem}

\begin{rem}
\source{tates-thesis-de-rham-setting}
\notready

In the physicist's language, we are 
comparing line operators for the $A$-twist of a pure hypermultiplet
with the $B$-twist of a $U(1)$-gauged hypermultiplet. Moreover, we
consider these $3d$ $\sN = 4$ theories as boundary theories for 
(suitably twisted) pure $U(1)$-gauge theory (with electro-magnetic duality
exchanging its $A$ and $B$-twists).

\end{rem}

\subsection{Connections to $3d$ mirror symmetry}

The equivalence of Theorem \ref{t:main} is a first\footnote{Specifically, 
as an equivalence of categories of line operators, with both matter and
a non-trivial gauge group on the $B$-side.}
instance of what is expected to be a broad family of equivalences, 
which goes under the heading
\emph{$3d$ mirror symmetry}. 
This subject is closely tied to the geometric Langlands program.

As context for our results, we
provide an informal introduction to these ideas below.  
Our objective is to connect our equivalence as closely as possible with the 
mathematical physics literature, and to promote those ideas to mathematicians
interested in the circle of ideas around geometric Langlands.

We hope the reader will forgive the 
redundancy of our discussion given 
the numerous other great sources in the literature.
% other refs?

We emphasize that we do not claim originality for any of the ideas
appearing below. We withhold attributions and leave much contact
with the existing literature until \S \ref{ss:3d-refs}.

\begin{rem}
\source{tates-thesis-de-rham-setting}
\notready

For another exposition of this subject also targeted at mathematicians,
see \cite{bf-notes}. For a recent physics-oriented exposition, we refer
to \cite{dggh} \S 2.

\end{rem}

\subsubsection{}

The physics of the last thirty years has highlighted the role of \emph{dualities}:
quantum field theories that are equivalent by non-obvious means.
We emphasize that
physicists regard these dualities as conjectural: they do not claim
to know how to match the QFTs, only (at best) parts of them.

Examples abound. For instance,
Montonen--Olive's (conjectural) $S$-duality for 
$4$-dimensional Yang-Mills theory with $\sN = 4$ supersymmetry for two 
Langlands dual gauge groups $G$ and $\check{G}$ is of keen interest. 

The dualities physicists consider fit into a sophisticated logical hierarchy.
But briefly, most are subsumed in the existence of Witten's (conjectural) $\sM$-theory,
which is supposed to recover other (known) QFTs by various constructions.

\subsubsection{}

A general problem is to extract mathematical structures
from quantum field theories.\footnote{Often, this is by a sort of analogy. Physicists
are drawn to differential geometry: their bread and butter are moduli spaces
of solutions to non-linear PDEs (namely, Euler-Lagrange equations).

But for a mathematician,
it may be preferable to consider algebraic varieties as analogous to K\"ahler
manifolds, or symplectic varieties as analogous to hyper-K\"ahler manifolds. 
Often, rich algebraic geometry results.} 
In this case, physical dualities relate to mathematical
conjectures.

For instance, in \cite{kapustin-witten}, 
Kapustin and Witten found such a relationship 
between Montonen--Olive's (conjectural) duality for 
$4$-dimensional Yang-Mills with $\sN = 4$ supersymmetry and geometric Langlands
conjectures in mathematics. Their perspective was developed in 
\cite{costello-ss-2-4} and 
\cite{elliott-yoo}, which exactly clarified some means of extracting
algebro-geometric invariants from Lagrangian field theories.

\subsubsection{$3d$ $\sigma$-models}

In one setting, so-called\footnote{We remind what these parameters
encode in the subsequent discussion.}
$3d$ $\sN = 4$ quantum field theories, there are 
rich relations with algebraic geometry.

First, for every algebraic stack $\sX$, we suppose
there is (in some algebraic sense) 
a $3d$ (Lagrangian) quantum field theory
$T_{\sX}$ with $\sN = 4$ supersymmetry; 
we spell out some more precise expectations below. 

\begin{rem}
\source{tates-thesis-de-rham-setting}
\notready\label{r:no-qft}
\source{tates-thesis-de-rham-setting}

At the quantum level, physicists agree our supposition is problematic: 
cf. the discussion at the
end of \cite{rozansky-witten} \S 2.3. (The 
classical field theory is fine, but may be unrenormalizable, and even if it is
renormalizable, there may not be a distinguished scheme.)
So it is better 
to regard $T_{\sX}$ as a heuristic:
its twists (see below) are all that is defined
(at the quantum level). 

\end{rem}

\begin{rem}
\source{tates-thesis-de-rham-setting}
\notready\label{r:3d-t*}
\source{tates-thesis-de-rham-setting}

In fact, physicists
say that $T_{\sX}$ depends only on the symplectic stack 
$T^* \sX$. At the classical level, this theory
is the $3d$ $\sigma$-model with target $T^*\sX$. 
The $\sN = 4$ supersymmetry comes from the 
symplectic structure on $T^*\sX$. 

I.e., more generally, there is a \emph{classical} $3d$ $\sN = 4$ theory
corresponding to a $\sigma$-model with target any symplectic stack $\sS$. 
One hopes for a quantization associated with a Lagrangian
foliation $\lambda:\sS \to \sX$; 
at the moment, this quantization (or rather, its $A/B$ twists) 
is only understood for $\sS$ a twisted cotangent bundle over some $\sX$
(and $\lambda$ the canonical projection).

This perspective is important, but sometimes 
inconvenient (especially at the quantum level), 
so we emphasize the role of $\sX$ e.g. in our notation.

\end{rem}

% reference for A&B twists specifically? A little muddled in [ES]....

\subsubsection{Algebraic QFTs}

Next, we describe how we think about a $3d$ QFT $T$ in algebraic geometry.
We include this discussion to make precise the connection between the mathematical
physics we wish to discuss and the algebraic geometry we wish to study.

\subsubsection{}

Recall from \cite{lurie-tft} that an \emph{oriented $3$-dimensional fully-extended TFT} 
$T$ would 
attach a number $T(M)$ to a closed 3-manifold $M$, a vector space
$T(M)$ to a closed 2-manifold $M$, a DG category $T(M)$ to a closed
1-manifold $M$, and a DG 2-category (i.e., $\DGCat_{cont}$-module
category) to a 0-manifold $M$. Cobordisms define morphisms. Disjoint unions
go to tensor products.

A variant: given an ambient $3$-manifold $N$, a fully-extended TFT \emph{on $N$} should
assign such data to manifolds $M$ equipped with embeddings into $N$.
(We are not aware of a reference.)

\subsubsection{}\label{ss:alg-3d}
\source{tates-thesis-de-rham-setting}

The algebraic situation is similar, but only some data is defined.
We heuristically describe the main structures, without giving formal definitions.

Fix a smooth, projective, geometrically connected algebraic curve $X$,
which we regard as analogous to a $2$-manifold. (The curve $X$ should not be confused
with the $\sX$ that sometimes occurs when considering a $\sigma$-model with 
target $T^*\sX$.)

A\footnote{As the input to the discussion is an algebraic curve $X$,
we sometimes refer to this formalism as \emph{algebraic QFT}.
We emphasize that we are \emph{not} referring to some kind of
``algebraic twist" (analogous to a ``holomorphic twist") 
of a supersymmetric Lagrangian field theory \textemdash{}  the
formalism is insensitive to such considerations. Rather, this
terminology refers to the fact that this formalism is defined without
analysis, and makes sense in the general setting of algebraic geometry (of
curves here, with evident adaptations in the terminology otherwise)
over fields (perhaps of characteristic zero).}
\emph{$3d$ field theory} $T$ on\footnote{It might be better
to say the theory lives on $X \times \bR$.} $X$ includes the
data of a vector space $T(X) \in \Vect$ and a category 
$T(\o{\sD}_x) \in \DGCat_{cont}$ for every point $x \in X$,
regarding $\o{\sD}_x$ as analogous to a circle. 

We regard the formal disc $\sD_x$ as a cobordism
$\emptyset \to \o{\sD}_x$. There is an associated functor:
%
\[
\Vect = T(\emptyset) \to T(\o{\sD}_x) \in \DGCat_{cont} 
\]

\noindent i.e., there is a preferred object of $T(\o{\sD}_x)$,
the so-called \emph{vacuum} (or \emph{unit}) object. 

We regard $X \setminus x$ as a cobordism
$\o{\sD}_x \to \emptyset$. There is an associated functor:
%
\[
T(\o{\sD}_x) \to T(\emptyset) = \Vect \in \DGCat_{cont}.
\]

\noindent This functional evaluated on the vacuum is
$T(X) \in \Vect$.

\begin{rem}
\source{tates-thesis-de-rham-setting}
\notready

Unlike in the topological situation, the above is not symmetric.
I.e., we only allow cobordisms in the directions specified above.

\end{rem}

\begin{rem}
\source{tates-thesis-de-rham-setting}
\notready

In this $3d$ setting, the category $T(\o{\sD}_x)$ is often
called the \emph{category of line operators} for the theory. (Physicists would 
draw the line $x \times \bR \subset X \times \bR$ passing through
the interior of our circle $\o{\sD}_x \times 0 \subset X \times \bR$.)

\end{rem}

\begin{rem}
\source{tates-thesis-de-rham-setting}
\notready\label{r:fact-cats}
\source{tates-thesis-de-rham-setting}

The assignment $x \rightsquigarrow T(\o{\sD}_x)$ should extend
to a \emph{factorization category on $X$} in the sense of
\cite{chiralcats}. The vacuum objects should correspond to a unital
structure on this factorization category. The functionals
above should extend to a well-defined functional on the
\emph{independent category} of this unital factorization category;
see \cite{indep-cat} for the definition.

\end{rem}

\begin{example}
\source{tates-thesis-de-rham-setting}
\notready\label{e:3d-chiral}
\source{tates-thesis-de-rham-setting}

In \cite{chiral}, Beilinson and Drinfeld defined such a structure
for any chiral algebra $\sA$ on $X$, i.e., they (in effect)
defined a $3d$ field theory $T_{\sA}$. 
The vector space $T_{\sA}(X)$ is the
space of \emph{conformal blocks} of $\sA$. The category
$T_{\sA}(\o{\sD}_x)$ is $\sA\mod_x$, the category of (unital) chiral
modules for $\sA$ supported at $x$. The vacuum object is the
vacuum representation of $\sA$, and the functional
$\sA\mod_x \to \Vect$ is the functor $C^{ch}(X,\sA,-)$ from
\emph{loc. cit}. \S 4.2.19.

(This theory is analogous to the $3d$ TFT associated
with a $\sE_2$-algebra $\sA$ in 
\cite{lurie-tft} Theorem 4.1.24.\footnote{We 
note that \emph{loc. cit}. regards this theory as a fully-extended
$2d$ theory with more categorical complexity than we outlined
before. We instead regard the theory as a not-fully-extended
$3d$ theory of the type outlined above. 
This compares to Remark 4.1.27 from \cite{lurie-tft}.})

\end{example}

\begin{rem}
\source{tates-thesis-de-rham-setting}
\notready

Informally, it is good to think of $3d$ theories on $X$ as 
the natural home for the Morita theory of chiral algebras,
just as $\DGCat_{cont}$ is the natural home for Morita theory of
associative algebras.

\end{rem}

\begin{rem}
\source{tates-thesis-de-rham-setting}
\notready\label{r:voa-3d}
\source{tates-thesis-de-rham-setting}

The relationship between chiral algebras and $3d$ $\sN = 4$ (and
sometimes $\sN = 2$) theories is the starting point of the series of
works \cite{voa-3d}, \cite{ccg}, and \cite{cdg}, which aim to use
chiral algebras to 
study mirror dual theories in this way. To describe this work in more detail,
we use ideas reviewed later in this introduction. 

Specifically, a $3d$ $\sN = 4$
theory $T$ with a suitably deformable 
(cf. \cite{voa-3d} \S 2.3; there is an implicit
choice of $A$ or $B$ twists fixed in our discussion here) 
$\sN = (0,4)$ boundary condition $\sB$ should yield an 
algebraic $3d$ field theory $T^{alg}$ 
and a boundary condition $\sB^{alg}$ for it; here
we regard $\sB^{alg}$ as an interface $T^{alg} \to \on{triv}$ 
\emph{to} the trivial theory. 
This data yields a factorization algebra $\sA_{T,\sB}$ on $X$: its fiber at $x$
is by evaluating $\sB^{alg}(\o{\sD}_x):T(\o{\sD}_x) \to \on{triv}(\o{\sD}_x) = \Vect$
at the vacuum object for $T$. In this case, there is a canonical
interface $T^{alg} \to T_{\sA_{T,\sB}}$. Sometimes, it is even an equivalence,
meaning that one can study the full theory $T^{alg}$ via the factorization
algebra $\sA_{T,\sB}$. 

\end{rem}

\begin{rem}
\source{tates-thesis-de-rham-setting}
\notready

As David Ben-Zvi emphasized to us, unital factorization categories
with a functional on their independent categories
do seem to provide a robust
definition of $[1,2]$-extended $3d$ algebraic QFT (while allowing for
some quite degenerate theories). He also informs us that
these ideas were developed in collaboration with Sakellaridis and Venkatesh,
and will be developed in their forthcoming work. 

\end{rem}

\begin{question}
\source{tates-thesis-de-rham-setting}
\notready

Is there some sense in which the theories of Example \ref{e:3d-chiral}
are $[0,2]$-extended, as for their $E_2$-counterparts? 
And what general constructions of morphisms
(alias: \emph{interfaces}) between such theories exist?\footnote{This question
is of practical importance to us. Our
main theorem amounts to a construction of a non-trivial interface 
between two $3d$ theories. It is desirable to understand this construction on 
better conceptual grounds.}

\end{question}

\subsubsection{Supersymmetry}

We now recall the main practical 
application of supersymmetry: supersymmetric
QFTs (typically) come with canonical 
deformations,\footnote{Sometimes, we think of
deformations \emph{as} quantizations. For us, the difference
is as follows.

Briefly, in $0$-dimensions, $\sP_0$-algebras can be quantized
by $\sE_0$-algebras, 
i.e., pointed vector spaces. On the other hand, pointed vector
spaces $V$ may be further \emph{deformed} by 
finding a filtered vector space $V^{def}$ with $\gr_{\dot} V^{def} = V$.
(In homological settings, such a structure is equivalent to equipping $V$ with 
a differential in a suitable sense, or one might say, deforming
the differential on $V$.)

In higher dimensions, one can say essentially the same words, following
\cite{costello-gwilliam-i}. Lagrangian
densities give rise to factorization $\sP_0$-algebras, 
which govern the corresponding
classical field theory. These may be quantized by factorization
$\sE_0$-algebras, i.e., factorization algebras; this amounts to 
quantizing the classical field theory (or more specifically, constructing
the OPE for local operators). These factorization algebras
may be further deformed to other factorization algebras.}
called \emph{twists}.\footnote{For instance,
for a manifold $X$, its de Rham complex deforms
its complex of global differential forms (considered
as equipped with zero differential). Similarly,
for $X$ a smooth
algebraic variety, its de Rham complex 
deforms its Hodge cohomology (or rather, the cochain
analogue). These examples appear in supersymmetric
quantum mechanics.} Twisting preserves dimensions
but lowers the amount of supersymmetry; in our examples,
there is no residual supersymmetry.

\begin{rem}
\source{tates-thesis-de-rham-setting}
\notready

See e.g. \cite{costello-ss-2-4} and \cite{costello-scheimbauer} 
for an introduction to the twisting procedure. See 
\cite{elliott-safronov} and \cite{esw}
for a detailed classification of the twists of supersymmetric QFTs. For the specific
twists considered in $3d$ mirror symmetry, see \cite{dggh} \S 2.2.

\end{rem}

\subsubsection{}

In the setting of $3d$ QFTs $T$ with $\sN = 4$ supersymmetry, 
there are two twists of interest for us: the $A$-twist $T_A$ and
the $B$-twist $T_B$. In practice, these are \emph{algebraic} theories.

For our theories $T_{\sX}$ of interest,
we let $T_{\sX,A}$ and $T_{\sX,B}$ denote the corresponding
$3d$ QFTs. 

\begin{example}
\source{tates-thesis-de-rham-setting}
\notready

The category $T_{\sX,A}(\o{\sD}_x)$ of line operators for the $A$-twisted theory
should be the category\footnote{At this point,
we are operating heuristically and do not want to be overly prescriptive, so
we do not consider finer points such as $D^!$ vs. $D^*$. Similarly in 
Example \ref{e:b}.

With that said, in our example, $D^!$ is what appears.}
$D(\fL\sX)$ of $D$-modules on the algebraic loop space: 
%
\[
\fL \fX = \Maps(\o{\sD}_x,\sX)
\]

\noindent of $\sX$.

\end{example}

\begin{example}
\source{tates-thesis-de-rham-setting}
\notready\label{e:b}
\source{tates-thesis-de-rham-setting}

The category $T_{\sX,B}(\o{\sD}_x)$ of line operators for the $B$-twisted
theory should be:
%
\[
\IndCoh\big(\Maps(\o{\sD}_{x,dR},\sX)\big).
\]

\end{example}

\begin{rem}
\source{tates-thesis-de-rham-setting}
\notready

If $\sX$ is smooth affine, each of the categories above essentially
come from chiral algebras, as in Example \ref{e:3d-chiral}.
Indeed, up to mild corrections, 
the $A$-twist is essentially governed by a CDO for $\sX$, while the
$B$-twist is governed by the commutative chiral algebra of functions 
on\footnote{As in \cite{chiral}, it is convenient to think of 
$\sX$ as $\Maps(\sD_{dR},\sX)$ here.} $\sX$.

\end{rem}

\begin{rem}
\source{tates-thesis-de-rham-setting}
\notready

The algebraic QFTs appearing above are of a special nature. First, they
are defined functorially on every smooth curve $X$. Moreover, for 
$X = \bA^1$, the resulting sheaves of categories are strongly
$\bG_a$-equivariant (cf. \cite{butson-thesis-i} Chapter 2).
These observations are a sort of de Rham analogue of the
fact that the $A$ and $B$-twists are \emph{topological} twists. 

\end{rem}

\subsubsection{} Below, we speak of both algebraic and non-algebraic (or \emph{analytic})
theories. Typically, we speak of supersymmetry for the non-algebraic theories, 
and obtain algebraic theories by twisting. 

\subsubsection{$3d$ mirror symmetry (first pass)}
%\hspace{-.15cm}\footnote{Specifically, the gauge theory in 
%question is pure gauge theory in $4$-dimensions, cf. \S \ref{sss:s} and the subsequent 
%discussion. (If $\sX$ is a stack but not a scheme, physicists might say 
%$T_{\sX}$ is a gauge theory; certainly they would when
%$\sX$ is equipped with a representable map to $\bB G$. Our emphasis here is not the 
%distinction between stacks and schemes, but rather 
%the additional symmetry of being a boundary
%condition for a $4d$ (pure) gauge theory.)} 

Given a $3d$ $\sN = 4$ theory $T$, there is another $3d$ $\sN = 4$ 
theory $T^{\star}$; this theory has the same underlying QFT as $T$, but the
\emph{embedding} of the supersymmetry algebra is modified by an automorphism
of that algebra. The salient property for us is that the $A$-twist 
$T_A^{\star}$ of $T^{\star}$ is the $B$-twist $T_B$ of the original theory.
Moreover, $(T^{\star})^{\star} = T$. We refer to $T^{\star}$ as the
\emph{abstract mirror dual} to the theory $T$.

\subsubsection{}

We now describe a simplified version of $3d$ mirror symmetry conjectures.

These conjectures state that for certain \emph{3d mirror dual pairs} 
$(\sX,\sX^{\star})$ of algebraic stacks, there is
an equivalence:
%
\[
T_{\sX}^{\star} \simeq T_{\sX^{\star}}
\]

\noindent of $3d$ field theories.

For instance, passing to $B$-twists and categories of line operators on both sides, 
such a conjecture predicts an equivalence:
%
\[
D(\fL\sX) \simeq \IndCoh\big(\Maps(\o{\sD}_{x,dR},\sX^{\star})\big).
\]

\begin{example}
\source{tates-thesis-de-rham-setting}
\notready\label{e:mirror-tate-ungauged}
\source{tates-thesis-de-rham-setting}

The pair $(\sX,\sX^{\star}) = (\bA^1,\bA^1/\bG_m)$ is supposed to be a $3d$ mirror
dual pair of stacks. In this case, Theorem \ref{t:main} amounts to 
an equivalence of line operators for the corresponding twisted\footnote{This
may seem incomplete, given that $3d$ mirror symmetry is stated more symmetrically 
in terms of untwisted theories. However, as in Remark \ref{r:no-qft}, 
the untwisted QFTs are on unsteady ground.}  
theories.

\end{example}

\subsubsection{$S$-duality}\label{sss:s}
\source{tates-thesis-de-rham-setting}

The theory of $3d$ mirror symmetry as presented above admits a generalization
in which there is an auxiliary pair $(G,\check{G})$ of Langlands dual
reductive groups; the previous setting amounts to a pair of trivial groups.

The theory becomes much richer in this setting. There are connections
to $S$-duality in $4d$ gauge theory and the geometric Langlands correspondence.
Moreover, many of the examples in conventional (i.e., trivial group) $3d$ mirror
symmetry can be better understood as being built from more primitive 
examples in the general setting. 

The major cost is that even the formulation of the conjectures 
becomes conditional on forms of $S$-duality conjectures.  

We discuss these ideas in more detail below. 

\subsubsection{}

First, we need to discuss $4d$ algebraic QFTs, which we denote $\mathsf{T}$.

This is easy: the formalism is the same as in \S \ref{ss:alg-3d},
but one categorical level higher. For instance, there should 
be a DG category attached to $X$, and a DG 2-category $\mathsf{T}(X)$
(of \emph{surface defects}) attached to
$\o{\sD}_x$.

We can speak of \emph{interfaces} $\mathsf{T}_1 \to \mathsf{T}_2$ between 
theories, which are morphisms in the natural sense. An interface 
$\mathsf{triv} \to \mathsf{T}$ from
the trivial\footnote{This theory attaches $\Vect$ to $X$ and 
$\DGCat_{cont}$ to $\o{\sD}_x$.}
theory is a \emph{boundary condition} for $\mathsf{T}$.
This amounts to objects:
%
\[
\begin{gathered}
\sB_X \in \mathsf{T}(X) (\in \DGCat_{cont})  \\
\sB_x \in \mathsf{T}(\o{\sD}_x) (\in \mathsf{2DGCat}).
\end{gathered}
\] 

\begin{rem}
\source{tates-thesis-de-rham-setting}
\notready\label{r:3d-as-bdry}
\source{tates-thesis-de-rham-setting}

A boundary condition $\mathsf{triv} \to \mathsf{triv}$ is the same
as a $3d$ theory. 

\end{rem}

\begin{rem}
\source{tates-thesis-de-rham-setting}
\notready

Note that theories can be tensored. The $4d$ theories we consider are naturally
self-dual. So we can (and do) consider interfaces and boundary conditions as
operating in both directions. In such a situation, given boundary conditions 
$\sB_1,\sB_2$ for $\mathsf{T}$, we let $\pair{\sB_1}{\sB_2}$ denote 
the resulting boundary condition $\mathsf{triv} \to \mathsf{triv}$
obtained by composing $\mathsf{triv} \xar{\sB_1} \mathsf{T} \xar{\sB_2} \mathsf{triv}$
and applying Remark \ref{r:3d-as-bdry}.

\end{rem}


\subsubsection{} 

We now discuss twists.

For $4d$ $\sN = 4$ (non-algebraic) theory $\mathsf{T}$, there are supposed
to be $A$ and $B$ twists $\mathsf{T}_A$ and $\mathsf{T}_B$.
Again, there is an involutive operation of 
\emph{abstract mirror dual} $\mathsf{T}^{\star}$, exchanging $A$ and $B$ twists.

%Given a $3d$ theory $T$, we may speak of it 

\subsubsection{}

We now recall that in $4d$, for a reductive group $G$, 
there is Yang-Mills $\YM_G$ theory with $\sN = 4$ 
supersymmetry. Again, there are $A$ and $B$-twists with algebraic
meaning.

\begin{example}
\source{tates-thesis-de-rham-setting}
\notready

The $A$-twisted theory $\YM_{G,A}$ attaches to $X$ the DG category 
$D(\Bun_G(X))$ of $D$-modules on $\Bun_G(X)$, and to $\o{\sD}_x$ 
the DG 2-category $\fL G\mod$ of DG categories with a (strong)
action of the loop group $\fL G$ to $\o{\sD}_x$.
For the latter, the vacuum object is $D(\Gr_G)$ and the (``chiral homology")
functional $\fL G\mod \to \Vect$ is given by tensoring over $D^*(\fL G)$ with 
$D^*(\Bun_G^{level,x})$.

\end{example}

\begin{example}
\source{tates-thesis-de-rham-setting}
\notready

The $B$-twisted theory $\YM_{G,B}$ attaches to $X$ the DG category 
$\QCoh(\LocSys_G(X))$ of quasi-coherent sheaves on 
$\LocSys_G(X)$, and to $\o{\sD}_x$ attaches
the DG 2-category $\ShvCat_{/\LocSys_{G}(\o{\sD}_x)}$.

\end{example}

\subsubsection{}\label{sss:bx-ym}
\source{tates-thesis-de-rham-setting}

Given a stack $\sX$ with a $G$-action, there is a corresponding
boundary condition $\sB_{\sX}$ for $\YM_G$.

For $\sX = G$, this boundary condition is called the \emph{Dirichlet}
boundary condition; we denote it by $\Dir_G$.

For $\sX = \Spec(k)$, this boundary condition is called the \emph{Neumann}
boundary condition; we denote it by $\Neu_G$.

For any boundary condition $\sB$ of $\YM_G$, we let $T_{\sB}$ denote
the $3d$ theory $\pair{\sB}{\Dir_G}$. By standard compatibility of $3d$ and $4d$ 
supersymmetry algebras, any such $T_{\sB}$ has $\sN = 4$ supersymmetry.

For a $3d$ $\sN = 4$ theory $T$, the data of \emph{$G$-flavor symmetry} 
is the data of a boundary condition $\sB$ for $\YM_G$ and an identification
$T \simeq T_{\sB}$. Similarly, \emph{$G$-gauge symmetry} for $T$
is the data of a boundary condition $\sB$ and an identification
$T \simeq \pair{\sB}{\Neu_G}$.

The above is compatible with twists; $\sB$ gives a boundary condition 
$\sB_A$ for the $A$-twist $\YM_{G,A}$, and similarly for $B$-twists. 
The resulting $3d$ (algebraic) theories are $T_{\sB,A}$ and $T_{\sB,B}$.

\begin{rem}
\source{tates-thesis-de-rham-setting}
\notready

For $\sX$ as above, the boundary condition $\sB_{\sX,A}$ should define
the object $D(\fL \fX) \in \fL G\mod$, and the boundary condition 
$\sB_{\sX,B}$ should define $\IndCoh(\Maps(\o{\sD}_{x,dR},\sX/G)) \in 
\ShvCat_{/\LocSys_G(\o{\sD})}$. Globally, there should be objects
of $D(\Bun_G)$ and $\QCoh(\LocSys_G)$ as well. These objects 
are considered by Ben-Zvi, Sakellaridis, and Venkatesh, who 
term them \emph{period/$L$-sheaves}, interpreting 
them as sheaf-theoretic analogues of period integrals/$L$-function values
from harmonic analysis and number theory. They were also considered
(away from singular loci) in \cite{gaiotto-s} \S 9-10. 

\end{rem}

\begin{rem}
\source{tates-thesis-de-rham-setting}
\notready

Parallel to Remark \ref{r:3d-t*}, at least at the classical (non-quantum) level,
the boundary condition $\sB_{\sX}$ can be defined more generally for
Hamiltonian $G$-spaces, with $\mu:T^*\sX \to \fg^{\vee}$ inducing
$\sB_{\sX}$ (and again, something weaker than 
a suitable Lagrangian foliation should be needed to quantize).

\end{rem}

\subsubsection{}

In the above setting, \emph{$S$-duality} is supposed to be an equivalence:
%
\[
\YM_G^{\star} \simeq \YM_G
\]

\noindent of $4d$ theories. Passing to $B$-twists, we obtain:
%
\[
\YM_{G,A} \simeq \YM_{\check{G},B}.
\]

\noindent Up to smearing\footnote{Throughout this exposition of the general ideas, 
we give ourselves the freedom of ignoring these subtleties. 
In the body of this paper (specifically, \S \ref{s:spectral} and \S \ref{s:indcoh}), 
we treat these technical points in detail in our context. The non-abelian future of the
subject will of course similarly need to confront these points more seriously.} 
away homological algebra subtleties 
(cf. \cite{arinkin-gaitsgory}), this is a form of the geometric Langlands
conjectures, encoding local and global theories and compatibilities
between them (cf. \cite{quantum-langlands-summary}). 

The subsequent discussion is predicated on the existence of $S$-duality, so 
we assume it in what follows. 

\subsubsection{}

Given a boundary condition $\sB$ for $\YM_G$, we let 
$\sB^{\star}$ denote the resulting \emph{abstract mirror dual}
boundary condition for $\YM_G^{\star}$, and then let
$\check{\sB}$ denote the resulting $S$-dual boundary condition for 
$\YM_{\check{G}}$.

\subsubsection{$3d$ mirror symmetry redux}

Now fix a reductive group $G$.

In general, a\footnote{Traditionally, \emph{3d mirror symmetry} refers to
the case where $G = \check{G} = \on{triv}$. We find the present terminology
convenient, although it departs somewhat from established conventions.}
\emph{$3d$ mirror dual pair (for $(G,\check{G})$)}
consists of a pair $(\sX,\sX^{\star})$ where $G \actson \sX$,
$\check{G} \actson \sX^{\star}$, and there is an equivalence:
%
\[
\check{\sB}_{\sX} \simeq \sB_{\sX^{\star}}
\]

\noindent of boundary conditions for $\YM_{\check{G}}$. I.e., the 
boundary condition $S$-dual to $\sB_{\sX}$ should be $\sB_{\sX^{\star}}$.
For $G$ trivial, this recovers our earlier notion of $3d$ mirror dual pairs.

\subsubsection{}

Suppose $(\sX,\sX^{\star})$ as above. Passing to $B$-twists, we find that the objects:
%
\[
D(\fL\sX) \in \fL G\mod \simeq 
\ShvCat_{/\LocSys_{\check{G}}(\o{\sD})} 
\ni 
\IndCoh(\Maps(\o{\sD}_{x,dR},\sX^{\star}/\check{G}))
\]

\noindent are supposed to match, where the middle isomorphism is local geometric
Langlands (considered modulo homological subtleties). 

We remark that this does not amount to an equivalence of categories
between the categories on the left and on the right here. Rather, this
statement is inherently conditional on local geometric Langlands. 
(And we intend it to be a little fuzzy regarding homological subtleties.) 

With that said, in the case that $G$ is \emph{abelian}, the above does
amount to an equivalence of categories compatible with certain symmetries. 
(In general, the Whittaker category for the LHS above should be close to the
RHS above.)

\subsubsection{Some examples of mirror dual pairs}

To illustrate the above, we briefly include some examples.
For many other examples, we refer to \cite{wang-table}, which is 
a table of examples maintained by Jonathan Wang. (The terminology in the
first two examples is taken from Ben-Zvi, Sakellaridis, and Venkatesh, 
who are appealing to the analogy with harmonic analysis.) 


\begin{example}[Godement-Jacquet]
\source{tates-thesis-de-rham-setting}
\notready\label{e:godement-jacquet}
\source{tates-thesis-de-rham-setting}

For $G = \check{G} = GL(V) \times GL(V)$, the pair
$(GL(V) \times V,\End(V))$ should be $3d$ mirror dual.

In particular, the objects:
%
\[
D(\fL GL(V) \times \fL V) \in \fL GL(V) \times \fL GL(V) \mod 
\]

\noindent should correspond to:
%
\[
\IndCoh(\Maps(\o{\sD}_{dR},GL(V) \backslash\End(V)/GL(V)) \in 
\ShvCat_{/\LocSys_{GL(V)\times GL(V)}(\o{\sD})}
\]

\noindent under local geometric Langlands. 

\end{example}

\begin{example}[Tate]
\source{tates-thesis-de-rham-setting}
\notready

For $V = \bA^1$, 
the above asserts that:
%
\[
D(\fL\bG_m \times \fL \bA^1) \in \fL \bG_m \times \fL \bG_m \mod 
\]

\noindent corresponds to:
%
\[
\IndCoh(\sY \times \LocSys_{\bG_m}).
\]

\noindent As our group is abelian, we expect an honest equivalence 
of categories. Passing to $\fL\bG_m$-invariants $\leftrightarrow$ 
$\{$fiber at $\on{triv} \in \LocSys_{\bG_m}\}$, we obtain a conjecture:
%
\[
D(\fL \bA^1) \simeq \IndCoh(\sY).
\]

\noindent This is our main result.  

\end{example}

\begin{example}[Gaiotto-Witten \cite{gaiotto-witten}, \S 3.3.1]
\source{tates-thesis-de-rham-setting}
\notready

Let $G = \check{G} = GL(V)$ for $\dim(V) = n$. 
Let ${\sL au}$ be the (Laumon) moduli space
of data $(V_1,V_2,\ldots,V_{n-1},f_1,\ldots,f_{n-1})$ where
$V_i$ is a vector space of dimension $i$ and $f_i:V_i \to V_{i+1}$ is a 
morphism, where by definition, $V_n = V$.

Then the pair:
%
\[
(\sX,\sX^{\star}) = (GL(V),{\sL au})
\] 

\noindent is expected to be $3d$ mirror dual (with respect to our $(G,\check{G})$).

Unwinding, this in effect gives a sort of formula (or kernel) for 
geometric Langlands for $GL_n$; unfortunately, the necessary symmetries on the
resulting object are not apparent.  

\end{example}

\subsubsection{Where do mirror dual pairs come from?}

We briefly address the stated question. There are two styles of answer.

First, one can try to reverse engineer examples. If one believes a mirror
dual $\sX^{\star}$ to some $\sX$ should exist, 
one can sometimes calculate the ring of
functions on $T^*{\sX^{\star}}$ as the Coulomb
branch\footnote{I.e., total cohomology of endomorphisms of the unit object
in $T_{\sX^{\star},B} = \IndCoh(\Maps(\o{\sD}_{dR},\sX))$.} of $\sX^{\star}$,
which must correspond to the Higgs branch\footnote{I.e., total cohomology 
of endomorphisms of the unit object
in $T_{\sX,A} = D(\fL \sX)$.} of $\sX$. 

Alternatively, as in the work of Ben-Zvi, Sakellaridis, and Venkatesh,
one can reverse engineer examples by fitting phenomena from number theory
into the above framework via the analogy between automorphic forms
and automorphic $D$-modules.

But sometimes (always?) there is a better answer than either. 
Many $3d$ mirror symmetry examples can be derived from (super)string 
theory dualities and the existence of $\sM$-theory. (For example, see 
\cite{gaiotto-witten} \S 3.1.2 for the derivation in our example.)
We are not aware of a mathematical counterpart to this idea, i.e., a simple conjectural
framework that subsumes all examples in $3d$ mirror symmetry. Clearly this
would be highly desirable. 

\begin{rem}
\source{tates-thesis-de-rham-setting}
\notready

Forthcoming work of the first author and Philsang Yoo will develop in detail the
mathematical derivation of mirror pairs from $S$-duality in string theory.

\end{rem}

\subsubsection{The role of this paper}

Our objective in writing this paper was to test the above ideas in the
simplest\footnote{Not covered by geometric Langlands/$S$-duality for $\YM$, 
and in which the full weight of the infinite dimensional geometry appears.} 
case of interest, which is the Tate case discussed above. 
Here $(G,\check{G}) = (\bG_m,\bG_m)$ and\footnote{The discrepancy here with the
ungauged version considered in Example \ref{e:mirror-tate-ungauged} is 
explained by the fact that $\Dir$ and $\Neu$ are $S$-dual for tori.} 
$(\sX,\sX^{\star}) = (\bA^1,\bA^1)$. 

As we consider abelian gauge groups,
for which geometric Langlands is unconditional, this example can be 
studied in complete detail. 
We perform this study at the level of categories of line operators. We will
return to global considerations (compatibility with chiral homology) in 
future work.

Given that we obtain complete positive results in this case, and
because many of the technical difficulties inherent in the $3d$ mirror symmetry
project occur in our example, we believe that our results provide strong
support for the general conjectures. 

\subsubsection{Some references}\label{ss:3d-refs}
\source{tates-thesis-de-rham-setting}

Many of the ideas discussed above were developed in 
collaborative work, often unpublished, of a number of mathematical physicists. 
We are grateful to many people who have shared these ideas
with us over the years, and find their intellectual generosity inspiring.
The downside of this situation is that we find some difficulty in accurately
attributing priority for the ideas. We do our best here, but apologize in 
advance for any omissions or inaccuracies, which are not intended as slights.

The first instances of $3d$ mirror symmetry were considered in 
\cite{intriligator-seiberg}. The connection with string theory dualities
was made in \cite{hanany-witten}. These ideas were developed further
in \cite{gaiotto-witten}, which considered interactions between the
\cite{hanany-witten} constructions and $S$-duality for super Yang-Mills.
In turn, \cite{gaiotto-s} translated those ideas into conjectures regarding
geometric Langlands via the Kapustin-Witten dictionary \cite{kapustin-witten}
between Langlands and $S$-duality. The physics literature on BPS line operators 
is too vast to survey here, but we refer the reader to \cite{dggh} for a
review of the literature in the $3d$ $\sN = 4$ context. 

The algebro-geometric interpretation of the $A$-side discussed
above is suggested by work of Braverman, Finkelberg and Nakajima
in \cite{nakajima-coulomb}, \cite{bfn}, \cite{bfn-slices}, 
which emphasized the role of Coulomb branches.
The description of categories of line operators in twisted $3d$ $\sN = 4$
theories was given in unpublished work of Kevin Costello, 
Tudor Dimofte, Davide Gaiotto, Philsang Yoo, and the first 
author.\footnote{These ideas were recorded in \S 1.1 of 
\cite{dggh} and in \cite{bf-notes} \S 7.} 
The earliest derivation
was based on applying the method of Elliott-Yoo \cite{elliott-yoo} 
in the $3d$ $\sN = 4$ context, 
i.e., calculating (derived) moduli spaces of solutions to Euler-Lagrange equations and 
quantizing (in the shifted symplectic sense).\footnote{Because 
we are not aware of a publicly available
account of this derivation, we direct the reader to \cite{costello-x-talks},
which is a series of talks Kevin Costello gave in 2014 that discuss some of 
the main ingredients. It is the earliest talk we know of that contains
these components.} 
By applying the yoga of $3d$ mirror symmetry, this description of 
categories of line operators also 
led to a conjectural form of Theorem \ref{t:main} 
(and other related examples).\footnote{For a physics-oriented discussion of this
work, see \cite{tudor-talk}.}
Philsang Yoo and the first author considered these line operator categories
in the framework of local geometric Langlands as a mathematical interpretation 
of \cite{gaiotto-witten} (inspired in part by \cite{bfn-ring}).\footnote{See
\cite{phil-talk} and \cite{jh-talk} for discussions of this work.}
On the $A$-side, an alternative derivation of the category of
line operators is suggested by 
\cite{voa-3d} \S 4; here the method is as described in Remark \ref{r:voa-3d}. 

Connections between $3d$ mirror symmetry 
and number theory (e.g., duality for spherical varieties and period
integrals \cite{sakellaridis-venkatesh}) were developed by Ben-Zvi, Sakellaridis, 
and Venkatesh, in forthcoming\footnote{The work is also highly publicized -- 
there are recorded lectures about the work readily available online 
(see e.g. \cite{bz-talk}).}
work, which has led to profound new insights and many new examples of 
dual pairs.

Finally, in addition to the above works, 
we have drawn inspiration particularly from \cite{bf-notes},
 \cite{blpw-hypertoric}, \cite{blpw}
\cite{bpw}, \cite{loop-spaces-and-connections},
\cite{char-field-thry}, 
\cite{krs}, \cite{ksv}, and \cite{teleman-icm} in our thinking about
mirror symmetry. 

\subsubsection{Some related works}

We also highlight some works that are closely related to ours.

First, \cite{ccg} studies a version of our problem
near the formal completion of $0 \in \sY$ (i.e., Theorem \ref{t:main} 
\emph{in perturbation theory}) using VOAs; in their picture (based on the
method of Remark \ref{r:voa-3d}), the
Lie algebra controlling the deformation theory is a 
$\mathfrak{psu}(1|1)$ Kac-Moody algebra. We understand that these ideas 
are currently being developed further by Andrew Ballin, Thomas Creutzig,
Tudor Dimofte, and Wenjun Niu.

Second, the works \cite{bfgt} and \cite{bft} 
of Braverman-Finkelberg-Ginzburg-Travkin also obtained geometric results from 
mirror symmetry conjectures, but in the non-abelian setting. 
There it is not possible at present
to prove that boundary conditions match under $S$-duality, since
local geometric Langlands is only conjectural for non-abelian $G$.
Therefore, the cited authors instead consider the categories obtained by 
pairing the vacuum boundary condition(s) for
$\YM_G$/$\YM_{\check{G}}$, and
establish the resulting conjectures for certain examples of 
$3d$ mirror dual pairs $(G,\check{G})$ and $(\sX,\sX^{\star})$. 

\subsection{Outline of the paper}

\subsubsection{} We now describe the contents of the present work.

\subsubsection{}\label{ss:intro-strategy}
\source{tates-thesis-de-rham-setting}

First, we describe the general strategy. 

By definition, there are fully faithful functors:
%
\[
D(\fL_n^+\bA^1) \into D^!(\fL^+ \bA^1) \into D^!(\fL \bA^1).
\]

\noindent Here $\fL_n^+\bA^1$ is the (scheme corresponding to) 
the vector space $k[[t]]/t^nk[[t]]$; $\fL^+ \bA^1$ is the
(scheme corresponding to the pro-finite dimensional) vector 
space $k[[t]]$, and $\fL\bA^1$ is the (indscheme corresponding to
the Tate) vector space $k((t))$. The first functor is a $!$-pullback
and the second functor is a (suitably normalized) $*$-pushforward functor.

The main step in our construction is the construction of corresponding
functors:
%
\[
\Delta_n:D(\fL_n^+\bA^1) \to \IndCoh^*(\sY).
\]

\noindent As the left hand side is modules over a Weyl algebra, this amounts
to constructing an object of the right hand side with an action of a Weyl
algebra. We construct this object explicitly, via generators and relations. 

From there, we bootstrap up to construct a functor
$\Delta$ as in Theorem \ref{t:main}. We show it is an equivalence using 
the geometry of $\sY$.

\begin{rem}
\source{tates-thesis-de-rham-setting}
\notready

There is an adage in geometric Langlands that most of the work occurs on the
geometric (i.e., automorphic) side. In our work it is the opposite: the 
geometric side is easy, and most of our work is on the spectral side. 

\end{rem}

\subsubsection{}

We introduce $\sY$ (and related spaces) in \S \ref{s:spectral}. Here we
give the main geometric tools in our study of $\sY$. We study
ind-coherent sheaves on $\sY$ in \S \ref{s:indcoh}, after some preliminary 
remarks about Abel-Jacobi maps in \S \ref{s:aj}. We construct the
functors $\Delta_n$ in \S \ref{s:weyl}; as indicated above, this is our main
construction. We check the corresponding compatibility of this construction 
with class field theory
in \S \ref{s:cft}. We then show fully faithfulness of $\Delta_n$ in 
\S \ref{s:ff}. We then prove Theorem \ref{t:main} in \S \ref{s:main}.

\subsection{Categorical conventions}

At various points, we use homotopical algebra and derived
algebraic geometry in Lurie's higher categorical form,
cf. \cite{htt}, \cite{higheralgebra}, \cite{sag}, though our notation more
closely follows the conventions of \cite{grbook}.

In general, our terminology should be assumed to be derived.
We refer to $(\infty,1)$-categories as \emph{categories} and 
$\infty$-groupoids as \emph{groupoids}. We let
$\Gpd$ denote the category of groupoids.

We let $\DGCat_{cont}$ denote the category of cocomplete\footnote{Rather, presentable.}
DG categories (over $k$) and continuous DG functors, cf. \cite{grbook} \S I.1.10.
We let $\otimes$ denote its standard (\emph{Lurie}) tensor product. 
We let $\Vect \in \DGCat_{cont}$ denote the DG category of (chain complexes
of) vector spaces. For a DG algebra $A/k$, we let $A\mod$ denote the
corresponding DG category. 
For a DG category $\sC$ and $\sF,\sG \in \sC$, we let
$\ul{\Hom}_{\sC}(\sF,\sG) \in \Vect$ denote the complex of
maps from $\sF$ to $\sG$. 

For a DG category $\sC$ with a $t$-structure, we use cohomological
indexing conventions and let $\sC^{\leq 0}$ denote the connective
objects, $\sC^{\geq 0}$ denote the coconnective objects, and
let $\sC^{\heart} \coloneqq \sC^{\leq 0} \cap \sC^{\geq 0}$ denote
the heart of the $t$-structure. Similarly, we let 
$\sC^+$ denote the eventually coconnective objects.


By default, \emph{schemes} are DG\footnote{It will turn out a posteriori 
that the primary objects
we consider are classical schemes (or stacks). But because we are studying
derived categories of coherent sheaves, the machinery and perspective of derived
algebraic geometry is quite convenient.}  
 schemes over $k$. We let
$\AffSch$ denote the category of affine schemes (over $k$),
and $\Sch$ the category of schemes (over $k$).

When we wish to refer to consider various non-derived objects as 
derived objects, we use the
term \emph{classical}. So we may speak of classical schemes, classical rings,
classical vector spaces, classical $A$-modules, etc., all of which 
are full subcategories of the corresponding (suitably) derived categories.

Finally, we sometimes use standard ideas and notation from the theory of
categories with $G$-actions. We refer to \cite{dario-*/!} for an introduction
to these ideas.  

\subsection{Acknowledgements}

We thank Sasha Beilinson, David Ben-Zvi, Dario Beraldo, 
Mat Bullimore, Dylan Butson, Sasha Braverman,
Justin Campbell, Kevin Costello, Gurbir Dhillon, Tudor Dimofte, Chris Elliott,
Tony Feng, Davide Gaiotto, Nik Garner, Dennis Gaitsgory, Sam Gunningham, 
Joel Kamnitzer, Ivan Mirkovic, Tom Nevins, Wenjun Niu, Pavel Safronov, 
Yiannis Sakellaridis, Ben Webster, Jonathan Wang, Alex Weekes, and
Philsang Yoo for many helpful and inspiring conversations related
to this work. The first author particularly wishes to acknowledge and thank 
Philsang Yoo for joint work conjecturing Theorem \ref{t:main}. 

J.H. is part of the Simons Collaboration on Homological Mirror Symmetry supported by 
Simons Grant 390287. S.R. was supported by NSF grant DMS-2101984.
This research was supported in part by Perimeter Institute for Theoretical Physics. 
Research at Perimeter Institute is supported by the Government of Canada through the 
Department of Innovation, Science and Economic Development Canada and by the Province of 
Ontario through the Ministry of Research, Innovation and Science.
